\section{Metode yang Direncanakan}\label{sec:metode}

\subsection{Kerangka Penelitian}\label{sec:kerangka}

Kerangka penelitian ditunjukkan Gambar~\ref{fig:kerangka}.

\begin{figure}[H]
  \centering
  \resizebox{\textwidth}{!}{% Kerangka penelitian. Kotak biru = variabel yang diuji.
\begin{tikzpicture}[
  box/.style={draw, rounded corners=2pt, align=center, font=\footnotesize,
              minimum height=1.15cm, minimum width=2.35cm, inner sep=2pt, fill=white},
  var/.style={box, fill=blue!8},
  arr/.style={-{Stealth[length=2mm]}, thick}]
\node[box] (img) at (0,0) {Citra JPG\\Tobacco3482};
\node[var] (ocr) at (3.0,1.35) {\textbf{Mesin OCR}\\Tesseract,\\EasyOCR, PaddleOCR};
\node[var] (emb) at (6.5,1.35) {\textbf{Embedding}\\FastText + CNN1D,\\BERT};
\node[box] (cnn) at (4.75,-1.35) {Cabang citra\\MobileNetV2\\(tetap)};
\node[var] (fus) at (9.6,0) {\textbf{Fusion}\\rata-rata logit,\\fitur + LN};
\node[box] (eva) at (12.6,0) {OA, macro F1,\\McNemar,\\$\Delta_\text{teks}$, oracle};
\node[font=\scriptsize, text=gray, align=center] (kotak) at (3.0,2.75)
     {kotak posisi teks $\rightarrow$ tahap lanjutan};
\draw[arr] (img.north) |- (ocr.west);
\draw[arr] (img.south) |- (cnn.west);
\draw[arr] (ocr) -- (emb);
\draw[arr] (emb.east) -| (fus.north);
\draw[arr] (cnn.east) -| (fus.south);
\draw[arr] (fus) -- (eva);
\draw[-{Stealth[length=1.6mm]}, gray, dashed] (ocr.north) -- (kotak.south);
\end{tikzpicture}
}
  \caption{Kerangka penelitian; kotak biru adalah variabel yang diuji}
  \label{fig:kerangka}
\end{figure}

Penelitian terdiri atas empat langkah:
\begin{enumerate}
  \item OCR ulang seluruh citra dengan tiga mesin.
  \item Melatih cabang teks untuk setiap kombinasi mesin OCR dan embedding.
  \item Menggabungkan setiap cabang teks dengan cabang citra yang sama.
  \item Mengevaluasi akurasi sekaligus pemanfaatan modalitas.
\end{enumerate}
Cabang citra sengaja dibuat tetap, supaya setiap perbedaan hasil hanya dapat berasal dari
sisi teks.

\subsection{Dataset dan Protokol}\label{sec:dataset}

\textbf{Tobacco3482} \parencite{kumar2014} berisi 3.482 citra dokumen dalam 10 kelas:
Advertisement, Email, Form, Letter, Memo, News, Note, Report, Resume, dan Scientific.
Kelas terbesar adalah Memo (620 dokumen) dan terkecil Resume (120 dokumen). Citra diambil
dari Kaggle dalam format JPG \parencite{kaggletobacco}, yang tidak memiliki lisensi baku.

\textbf{QS-OCR-Small} \parencite{qsocr} berlisensi Apache-2.0. Isinya teks Tesseract~4
dari citra TIF asli, dan dipakai sebagai \textbf{kondisi acuan}, yaitu teks yang sama
dengan yang dipakai \textcite{audebert2020}.

Protokol pembagian data mengikuti \textcite{audebert2020}:
\begin{itemize}
  \item 800 dokumen latih per pembagian, 10\% di antaranya dipakai untuk validasi;
  \item sisanya, 2.682 dokumen, menjadi data uji;
  \item tiga pembagian acak terstratifikasi, dengan seed 42, 43, dan 44;
  \item \textbf{indeks pembagian yang sama} dipakai di semua kondisi, supaya perbedaan
  antar-kondisi tidak tercampur perbedaan data.
\end{itemize}
Pemakaian dataset dibatasi untuk keperluan akademik non-komersial, dan tidak ada berkas
dataset yang didistribusikan ulang.

\subsection{Ekstraksi OCR Multi-mesin}\label{sec:metode-ocr}

Konfigurasi setiap mesin disajikan pada Tabel~\ref{tab:mesin}.

\begin{table}[H]
  \centering
  \caption{Mesin OCR dan konfigurasinya}
  \label{tab:mesin}
  \fontsize{10}{12}\selectfont
  \begin{tabularx}{\textwidth}{|l|>{\raggedright\arraybackslash}X|l|}
    \hline
    \textbf{Mesin} & \textbf{Konfigurasi} & \textbf{Kotak teks} \\ \hline
    QS-OCR (acuan) & Tesseract 4.0-beta pada TIF asli, \texttt{-{}-oem 1 -{}-psm 3} &
    tidak tersedia \\ \hline
    Tesseract & versi 5 pada JPG, konfigurasi sama dengan QS-OCR & per kata \\ \hline
    EasyOCR & model bahasa Inggris, GPU & per baris \\ \hline
    PaddleOCR & model PP-OCRv5 \textit{mobile}, CPU, tanpa koreksi orientasi &
    per baris \\ \hline
  \end{tabularx}
\end{table}

Empat keputusan desain menyertai konfigurasi tersebut:
\begin{enumerate}
  \item \textbf{Semua mesin membaca citra JPG yang sama}, termasuk Tesseract yang
  dijalankan ulang. Selisih antara QS-OCR dan Tesseract-JPG karena itu mengukur efek
  gabungan format citra dan versi Tesseract.
  \item \textbf{Urutan baca mesin berbasis kotak direkonstruksi}: baris dari atas ke
  bawah, kata dari kiri ke kanan. Karena CNN~1D peka urutan, perbedaan urutan baca dicatat
  sebagai bagian dari efek mesin.
  \item \textbf{Teks dan kotak posisi setiap halaman disimpan}, sehingga tahap lanjutan
  tidak memerlukan OCR ulang.
  \item \textbf{Versi pustaka dan nama model dikunci serta dicatat}, karena model bawaan
  mesin OCR dapat berganti antar-versi.
\end{enumerate}

Keras-OCR disarankan dosen, tetapi dimasukkan secara bersyarat. Alasannya dijelaskan pada
risiko R3 (Subbagian~\ref{sec:risiko}).

\subsection{Proksi Kualitas OCR}\label{sec:proksi}

Tobacco3482 tidak memiliki transkripsi acuan, sehingga \textit{character error rate}
tidak dapat dihitung. Tiga proksi berikut menggambarkan keluaran OCR tanpa menilai
benar-salahnya.

\begin{enumerate}
  \item \textbf{Kata per halaman dan persentase halaman kosong.} Proksi ini menangkap
  mesin yang melewatkan blok teks.
  \item \textbf{Laju OOV.} Untuk dokumen $d$, laju OOV adalah bagian kata alfabetik dengan
  panjang minimal empat huruf yang tidak terdapat dalam daftar kata bahasa Inggris $V$
  \parencite{wordfreq}:
  \begin{equation}\label{eq:oov}
    \mathrm{OOV}(d) = \frac{\left|\{w \in d : |w| \geq 4,\; w \notin V\}\right|}
                           {\left|\{w \in d : |w| \geq 4\}\right|}
  \end{equation}
  Salah baca yang kebetulan berupa kata yang ada, misalnya \texttt{tbe} untuk
  \texttt{the}, tidak terhitung. Karena itu nilai ini adalah \emph{batas bawah} galat.
  \item \textbf{Kesepakatan antar-mesin.} Kesepakatan dua mesin diukur dengan indeks
  Jaccard atas himpunan kata ($A$ dan $B$) yang keduanya baca:
  \begin{equation}\label{eq:jaccard}
    J(A, B) = \frac{|A \cap B|}{|A \cup B|}
  \end{equation}
  Urutan kata sengaja diabaikan, karena mesin berbeda urutan baca walaupun kata yang
  dibaca sama.
\end{enumerate}

\subsection{Representasi Teks dan Citra}\label{sec:representasi}

Dua representasi teks diuji.

\textbf{FastText dan CNN~1D}, mengikuti \textcite{audebert2020}. Setiap kata dipetakan ke
vektor FastText 300 dimensi (Persamaan~\ref{eq:fasttext}). Sekuens dibatasi 500 kata,
lalu diproses CNN~1D.

\textbf{BERT-base}, mengikuti \textcite{kashyap2026}. Model ini di-\textit{fine-tune}
dengan masukan maksimum 512 token. Pelatihan memakai \textit{mixed precision} dan
akumulasi gradien supaya muat di GPU 4~GB. Bila tetap tidak muat, model diganti
DistilBERT, dan penggantian itu dicatat sebagai deviasi.

Teks masuk ke kedua representasi \textbf{tanpa} praproses gaya Kashyap, yaitu tanpa
membuang angka dan \textit{stopword}. Tujuannya agar efek embedding tidak tercampur efek
praproses.

Cabang citra memakai MobileNetV2 berbobot awal ImageNet, dengan masukan $384\times384$
dan rasio aspek di-\textit{warp}, mengikuti \textcite{audebert2020}. Cabang ini sama di
semua kondisi.

\subsection{Strategi Fusion yang Diuji}\label{sec:fusion}

Dua strategi fusion diuji.

\textbf{Rata-rata logit} (Persamaan~\ref{eq:logit}), yaitu aturan
\textcite{kashyap2026}. Strategi ini tidak memerlukan pelatihan ulang, sehingga dapat
diterapkan pada \emph{semua} kondisi.

\textbf{Fusion tingkat fitur dengan normalisasi per cabang}:
\begin{equation}\label{eq:fusionfix}
  \hat{y} = \arg\max_c \; \mathrm{MLP}\!\left([\,\mathrm{LN}(\mathbf{v}_t) \,;\,
  \mathrm{LN}(\mathbf{v}_v)\,]\right)_c
\end{equation}
Layer Normalization (Persamaan~\ref{eq:layernorm}) menyamakan skala kedua vektor fitur
sebelum digabung, supaya head klasifikasi tidak otomatis didominasi cabang yang skala
fiturnya lebih besar. Cabang teks diinisialisasi dari model teks yang sudah dilatih
sendiri, supaya kedua cabang mulai dari titik yang sama-sama terlatih.

Fusion tingkat fitur memerlukan pelatihan \textit{end-to-end} yang relatif mahal, sehingga
hanya diterapkan pada kondisi OCR terbaik dan terburuk.

\subsection{Desain Eksperimen dan Hipotesis}\label{sec:desain}

Matriks eksperimen disajikan pada Tabel~\ref{tab:desain}.

\begin{table}[H]
  \centering
  \caption{Matriks eksperimen}
  \label{tab:desain}
  \fontsize{10}{12}\selectfont
  \begin{tabularx}{\textwidth}{|l|>{\raggedright\arraybackslash}X|c|}
    \hline
    \textbf{Kelompok} & \textbf{Kondisi} & \textbf{Jumlah run} \\ \hline
    Cabang teks & 4 kondisi OCR $\times$ 2 embedding $\times$ 3 seed & 24 \\ \hline
    Cabang citra & MobileNetV2 $\times$ 3 seed & 3 \\ \hline
    Rata-rata logit & setiap model teks + cabang citra dengan seed yang sama &
    24 (tanpa pelatihan) \\ \hline
    Fusion fitur & kondisi OCR terbaik dan terburuk, FastText, seed 42 &
    2 (+4 bila waktu cukup) \\ \hline
    Uji \textit{missing modality} & setiap model fusion & tanpa pelatihan \\ \hline
  \end{tabularx}
\end{table}

\textbf{Pemeriksaan kewarasan.} Kondisi QS-OCR $\times$ FastText harus menghasilkan
akurasi cabang teks yang dekat dengan angka \textcite{audebert2020} (0,738). Bila jauh
berbeda, ada yang salah dalam pipeline, dan eksperimen dihentikan sampai penyebabnya
ditemukan.

Hipotesis dirumuskan \emph{sebelum} eksekusi, dan akan dilaporkan apa adanya, termasuk
bila keliru.

\begin{enumerate}[label=H\arabic*.,leftmargin=1.2cm,labelwidth=1.2cm]
  \item \textbf{(RM1)} Perbedaan akurasi cabang teks antar-mesin OCR bersifat nyata,
  yaitu melampaui variasi antar-seed. Mesin modern belum tentu unggul. Tobacco3482 berisi
  pindaian cetak hitam putih, yaitu kondisi tempat Tesseract paling kuat. Selain itu,
  CNN~1D peka urutan baca, yang berbeda antar-mesin.
  \item \textbf{(RM2)} BERT memperkecil selisih antar-mesin dibandingkan FastText.
  Argumen tandingannya: model bahasa standar juga melemah pada teks OCR
  \parencite{michail2025}, dan batas 512 token memotong dokumen panjang.
  \item \textbf{(RM3)} Pada fusion fitur, $\Delta_\text{teks}$ membesar seiring akurasi
  cabang teks. Rata-rata logit menjamin $\Delta_\text{teks} > 0$ secara konstruksi,
  tetapi akurasinya dapat turun bila cabang teks lemah.
\end{enumerate}

\subsection{Evaluasi dan Uji Statistik}\label{sec:evaluasi}

Setiap model dievaluasi dengan empat ukuran: OA, macro F1
(Persamaan~\ref{eq:macrof1}), F1 per kelas, dan oracle. Untuk setiap model fusion
ditambahkan $\Delta_\text{teks}$ dan $\Delta_\text{citra}$ (Persamaan~\ref{eq:delta}).

Perbandingan dua kondisi memakai uji McNemar (Persamaan~\ref{eq:mcnemar}) pada data uji
yang sama, per seed. Uji binomial eksak dipakai bila $b + c < 25$. Ada enam pasangan
kondisi OCR per jenis embedding, sehingga ambang signifikansinya dikoreksi menjadi
$0{,}05/6$. Sebuah perbedaan dinyatakan konsisten bila signifikan dan arahnya sama di
ketiga seed. Semua angka dilaporkan sebagai rata-rata $\pm$ simpangan baku tiga seed.
